% Independently audited by root and four_regular; finite checks are recorded below.
% Include after random-hosts-phase3.tex and theta-random-host-partition.tex.
\section{A wheel construction and an upper coefficient of ten ninths}
\label{sec:rhost-ten-ninths}

This section adds one external matching theorem to the previous construction.
All metric distributions and random-host embedding estimates are proved below.
The matching input is Pippenger's near-perfect matching theorem.
We use its version allowing $o(v)$ exceptional vertices, as stated after Theorem~1.1 in
Alon and Yuster, \emph{On a hypergraph matching problem},
\url{https://web.math.princeton.edu/~nalon/PDFS/ayhyper2.pdf}.

\begin{lemma}[Wheel length packing]\label{lem:rhost-wheel-lengths}
For integers $h\to\infty$, there are $(1+o(1))h/12$ traceable subdivisions of the four-rim wheel.
Their cycle lengths cover all but $o(h)$ integers in $[h+1,2h]$.
Each subdivision has $O(h)$ edges, shortest cycle length at least $2\lceil\sqrt h\rceil$,
and at least $(r-1)h$ edges in every connected core of cycle rank $r\in\{2,3,4\}$.
\end{lemma}

\begin{proof}
Put $s=\lceil\sqrt h\rceil$.
Choose $U,L,M$ independently and uniformly from the integers $s,\ldots,2s$.
Put $t=h+U$ and
\[
 c_j=4^{-j}\binom{2j}{j}\qquad(j\ge0).
\]
Choose four nonnegative integers $X_0,\ldots,X_3$ of sum $t$, with probability
\begin{equation}
 \Pr(X_0=x_0,\ldots,X_3=x_3)
   =\frac{c_{x_0}c_{x_1}c_{x_2}c_{x_3}}{t+1}.
 \label{eq:rhost-wheel-distribution}
\end{equation}
This is a probability distribution because $\sum_{j\ge0}c_jz^j=(1-z)^{-1/2}$.
For every pair $i,j$, the sum $X_i+X_j$ is uniform on $\{0,\ldots,t\}$.
Indeed, each convolution of two sequences $(c_j)$ is identically one.

Give the four spokes lengths $w_i=h+s-X_i$.
Give the cyclic rim arms $01,12,23,30$ lengths $1,L,1,M$.
Discard the outcome if a spoke length is nonpositive.
The probability of this event is $o(1)$.
To see this, fix a spoke and put $r=t-X_i$.
The generating function for the other three coordinates gives
\[
 \Pr(t-X_i=r)=\frac{c_{t-r}(2r+1)c_r}{t+1}.
\]
Since $c_j\le C/\sqrt{j+1}$, summing over $r\le s$ gives $O((s/h)^{3/2})$.
A union bound handles all four spokes.

There are twelve non-rim cycles, two for each pair of spokes.
Their lengths have the form
\[
 Z_{ij,a}=2h+2s-X_i-X_j+a.
\]
The two possible offsets for each pair are
\[
\begin{array}{c|c}
 ij&\text{offsets}\\\hline
 01,23&1,\ L+M+1\\
 12&L,\ M+2\\
 30&M,\ L+2\\
 02,13&L+1,\ M+1.
\end{array}
\]
All these lengths lie in
$J=[h+1,2h+6s+1]\cap\mathbb Z$.
For every $z\in I=[h+5s+1,2h]\cap\mathbb Z$ and each specified cycle label,
\begin{equation}
 \Pr(Z_{ij,a}=z)=\mathbb E\frac1{h+U+1}
                  =\frac{1+O(s/h)}h.
 \label{eq:rhost-wheel-marginal}
\end{equation}
For every integer $z$, the same probability is at most $1/h$.

We next bound the joint probabilities of two distinct cycle labels.
For the same spoke pair, the difference of the offsets is one of
$\pm(L+M)$, $\pm(L-M)$, or $\pm(L-M+2)$.
For fixed target lengths, this relation has probability $O(1/s)$.
Conditioning on $U,L,M$ then costs at most $1/h$ for the required pair sum.
For complementary spoke pairs, the sum of the two lengths is
\[
 3h+4s-U+a+a'.
\]
Conditioning on $L,M$ fixes at most one value of $U$.
The remaining pair-sum condition again costs at most $1/h$.
Both cases therefore give $O(1/(hs))$.

Suppose that the spoke pairs share one index.
Condition on $X_i+X_j=q$.
The split within this pair and the split within its complement are independent.
A split of an integer $q$ has masses $c_a c_{q-a}$, whose maximum is at most $C/\sqrt{q+1}$.
The shared-index pair sum is a convolution of splits of $q$ and $t-q$.
Its largest mass is at most
\[
 \min\left\{\frac C{\sqrt{q+1}},\frac C{\sqrt{t-q+1}}\right\}
 \le\frac C{\sqrt h}.
\]
The first pair sum has mass $1/(t+1)$.
Thus, uniformly in two target lengths and two distinct labels,
\begin{equation}
 \Pr(Z=z,Z'=z')=O(h^{-3/2}).
 \label{eq:rhost-wheel-codegree}
\end{equation}
Here we used $s=\lceil\sqrt h\rceil$.
Summing over $z=z'\in J$ shows that repeated cycle lengths occur with probability $O(h^{-1/2})$.
Discard these outcomes too.

For each twelve-element subset $A\subset J$, let $p_A$ be the probability that the retained outcome has length set $A$.
Then
\[
 \sum_Ap_A=1-o(1),\qquad
 \sum_{A\ni z}p_A\le\frac{12}h,\qquad
 \max_{z\ne z'}\sum_{A\supset\{z,z'\}}p_A=O(h^{-3/2}).
\]
Before the discarded outcomes, every vertex of $I$ has labelled incidence weight $(12+o(1))/h$.
Discarding outcomes removes total incidence weight $o(1)$.
Consequently, all but $o(h)$ vertices of $J$ have degree weight $(12+o(1))/h$.

This weighted system reduces to a simple hypergraph.
Set $\lambda=h(\ln h)^2$ and include each distinct set $A$ independently with probability $\lambda p_A$.
These probabilities are below one for large $h$, by the codegree bound.
Chernoff bounds show that the maximum degree is at most $(12+o(1))(\ln h)^2$.
All but $o(h)$ vertices have degree $(12+o(1))(\ln h)^2$.
The maximum codegree is $O(\ln h)$.
For completeness, each degree is a sum of independent indicators with mean at most $12(\ln h)^2$.
Use an error $(\ln h)^{-1/4}$ and a union bound over $O(h)$ vertices.
Each codegree has mean $O(h^{-1/2}(\ln h)^2)$.
The upper-tail estimate at $C\ln h$ permits a union bound over $O(h^2)$ pairs.
Pippenger's theorem, with its exceptional-vertex extension, gives a matching covering $|J|-o(h)$ vertices.
It has $h/12+o(h)$ edges.
Choose one retained wheel metric for each matching edge.
This proves the length coverage.

We verify the graph conditions.
The total spoke length is $3h+4s-U$, and the total edge count is at least $3h+4s+2$.
Every spoke has length at most $h+s$.
Every pair of spokes has total length at least $h$.
A non-rim cycle therefore has length at least $h+1$, while the rim has length $L+M+2\ge2s$.

A connected core of rank two contains either two spokes and the full rim, or at least three spokes.
Its edge count is at least $h$.
A core of rank three contains either three spokes and the full rim, or all four spokes.
Its edge count is at least $2h$.
A core of rank four is the whole wheel and has at least $3h$ edges.
These classifications follow by writing the rank as $k+r-v$,
where $k$ counts spokes and $r,v$ count rim edges and rim vertices.

Finally, delete the unit rim edges $01,23$.
Delete the hub-end edge of spokes $0,1$.
The remaining graph is a spanning path, with branch order $1,2,\text{hub},3,0$.
If a deleted spoke has length one, its rim vertex is an endpoint.
This proves traceability.
\end{proof}

\begin{lemma}[Wheel embeddings]\label{lem:rhost-wheel-embeddings}
Let $d\sim\ln n/\ln\ln n$, $h=d+O(1)$, and $D=\Theta(d)$.
Suppose $nD^{-d}$ tends to zero faster than every fixed negative power of $d$.
Any list of $O(d)$ wheel subdivisions supplied by Lemma~\ref{lem:rhost-wheel-lengths}
embeds vertex-disjointly in $G(n,D/n)$ with high probability.
The embeddings can avoid any set of $O(d^2)$ vertices chosen after observing the graph.
\end{lemma}

\begin{proof}
The overlap proof of Lemma~\ref{lem:rhost-k4-moments} extends to this fixed wheel kernel.
Here are the required details.
A wheel has no two vertex-disjoint cycles, so an overlap has at most one cyclic component.
After removing leaves, its core consists of complete kernel arms.
There are only constantly many core types.
For a core of rank $r$, its normalized overlap weight is at most $n^{r-1}D^{-e_r}$.
The metric lemma gives
\[
 e_1\ge2\lceil\sqrt h\rceil,\quad e_2\ge h,\quad e_3\ge2h,\quad e_4\ge3h.
\]
On eight arms, at most sixteen cut positions specify a connected overlap.
There are at most sixteen leaves and five other vertices of degree different from two.
Choose their images and then the images of the suppressed path segments.
A graph of cycle rank four has at most sixteen simple paths between prescribed endpoints.
It has at most fifteen simple cycles, so loop segments also have only constantly many choices.
Thus the number of cyclic overlap descriptions is at most $C(v+1)^{48}$.
Forest overlaps contribute $O(v^2/n)$, as before; maximum degree four only changes the geometric-series constant.
For large $n$, $D$ exceeds that constant.

Both the relative variance and the normalized intersecting-pair expectation are at most
\[
 C\left[\frac{v^2}n+(v+1)^{48}
 \left(D^{-2\lceil\sqrt h\rceil}+nD^{-h}+n^2D^{-2h}+n^3D^{-3h}\right)\right].
\]
This tends to zero faster than every fixed negative power of $d$.
The Chebyshev, Markov, and independent-set argument from Lemma~\ref{lem:rhost-theta-list}
therefore gives more than $C_0d^2$ disjoint copies of each metric, for every fixed $C_0$.
A greedy choice proves the assertion, including the adaptive forbidden set.
\end{proof}

\begin{theorem}[An upper coefficient of ten ninths]\label{thm:rhost-ten-ninths}
For the same $p_*$ as in Theorem~\ref{thm:rhost-scale} and every $p\ge p_*$,
with high probability
\[
 \operatorname{Pex}(G(n,p))
 \le\log_2n+\left(\frac{10}9+o(1)\right)\frac{\ln n}{\ln\ln n}.
\]
\end{theorem}

\begin{proof}
Use the exposure scheme of Theorem~\ref{thm:rhost-eight-sevenths}.
Put $h=9\lceil d/9\rceil$.
Use single cycles through length $h/3+1$.
Use the constant-sum theta partition in Lemma~\ref{lem:rhost-theta-middle} for $[h/3+2,h-2]$.
Use Lemma~\ref{lem:rhost-wheel-lengths} for $[h+1,2h]$.
Add one single-cycle block for every still-uncovered length through $T=2(d+b+1)$.
There are $o(d)$ such extra lengths, since $T-2h=O(b)=o(d)$.
The rank sum of the short blocks is at most
\[
 \frac h3+\frac{4h}9+4\frac h{12}+o(d)
 =\left(\frac{10}9+o(1)\right)d.
\]
All these choices depend only on $n$, before exposing the random host.

Prepare the single cycles and the previous base cycles in $G_1$.
Embed the theta and wheel blocks in the first half of $G_2$.
The theta moment estimate applies because every theta cycle has length at least $h/3+2$ and its edge count is $h\ge d$.
Lemma~\ref{lem:rhost-wheel-embeddings} supplies the wheel blocks.
Each type has more than $C_0d^2$ disjoint copies, so both lists can be selected greedily together.
Every block has a spanning path and the total number of vertices is $O(d^2)$.
Use the second half of $G_2$ for the connector chain.
The remaining exposures and interval coverage are unchanged.
The final excess is
\[
 K+2b+m+\left(\frac{10}9+o(1)\right)d
 =\log_2n+\left(\frac{10}9+o(1)\right)\frac{\ln n}{\ln\ln n}.
\]
\end{proof}

The finite checker \texttt{check\_random\_host\_wheels.py} verifies 3962 wheel metrics, all 42 kernel core types, and 120 explicit spanning paths.
It also verifies 96 pair-marginal identities with exact rational arithmetic.
These checks do not certify the asymptotic matching or embedding arguments.
